% ============================================================
% Paper 45: Ring Theory and Commutative Algebra
% Author: Michael Fuhrmann
% Project: Specialist Mathematics, Build 135
% Date: 2026-03-12
% ============================================================

\documentclass[12pt,a4paper]{amsart}
\usepackage{amsmath,amssymb,amsthm}
\usepackage{mathtools}
\usepackage[utf8]{inputenc}
\usepackage[T1]{fontenc}
\usepackage{hyperref}
\usepackage{enumitem,booktabs,array}
\usepackage{geometry}
\geometry{margin=2.5cm}

% ---- Theorem environments ----
\newtheorem{theorem}{Theorem}[section]
\newtheorem{lemma}[theorem]{Lemma}
\newtheorem{corollary}[theorem]{Corollary}
\newtheorem{proposition}[theorem]{Proposition}
\newtheorem{conjecture}[theorem]{Conjecture}
\theoremstyle{definition}
\newtheorem{definition}[theorem]{Definition}
\newtheorem{example}[theorem]{Example}
\newtheorem{remark}[theorem]{Remark}

% ---- Operators ----
\DeclareMathOperator{\Spec}{Spec}
\DeclareMathOperator{\Ann}{Ann}
\DeclareMathOperator{\rad}{rad}
\DeclareMathOperator{\Nil}{Nil}
\DeclareMathOperator{\Jac}{Jac}
\DeclareMathOperator{\lcm}{lcm}
\DeclareMathOperator{\ht}{ht}
\DeclareMathOperator{\codim}{codim}
\DeclareMathOperator{\tr}{tr}
\DeclareMathOperator{\chr}{char}
\DeclareMathOperator{\Frac}{Frac}
\DeclareMathOperator{\Cl}{Cl}

\title{Ring Theory and Commutative Algebra\\
\large Ideals, Factorization, Noetherian Rings, and Gröbner Bases}

\author{Michael Fuhrmann}
\address{Specialist Mathematics Project, Build 135}
\email{specialist-maths@localhost}
\date{\today}

\begin{document}

\begin{abstract}
This paper presents a self-contained introduction to ring theory and commutative
algebra, covering the central objects and theorems of modern algebra.
Starting from the axiomatic definition of a ring, we develop the theory of ideals
and quotient rings, establish the isomorphism theorems, and study polynomial rings
with their irreducibility criteria. The hierarchy of special ring classes---Euclidean
domains, principal ideal domains (PIDs), and unique factorization domains (UFDs)---is
carefully developed. We then treat Noetherian and Artinian rings, including the
Hilbert basis theorem and Nakayama's lemma. Localization and its application to
Dedekind rings provide the bridge to algebraic number theory. The paper concludes
with an introduction to Gröbner bases and the Buchberger algorithm as a
computational tool. All results are accompanied by proofs or proof sketches and
are illustrated with concrete examples. The exposition is based on the
implementations in \texttt{ring\_theory.py} and \texttt{commutative\_algebra.py}
of the Specialist Mathematics project.
\end{abstract}

\maketitle
\tableofcontents

% ============================================================
\section{Introduction}
% ============================================================

Rings are the fundamental algebraic structures that generalize the arithmetic of
integers. Wherever addition and multiplication interact---in number theory, geometry,
combinatorics, or mathematical physics---ring-theoretic methods provide the
unifying framework.

The historical development of ring theory was driven by concrete problems.
Gauss studied $\mathbb{Z}[i]$ to prove his two-squares theorem. Kummer introduced
``ideal numbers'' to salvage unique factorization in cyclotomic fields. Dedekind
formalized the modern notion of an ideal and proved that rings of algebraic integers
are what we now call Dedekind domains---where unique factorization holds at the
level of ideals even when it fails for elements.

In the twentieth century, Emmy Noether~\cite{Noether1921} revolutionized the subject
by shifting focus from explicit calculations to structural principles. Her ascending
chain condition, the Noetherian property, turned out to be the exact hypothesis
needed to make algebra work cleanly. At the same time, Hilbert's basis theorem
showed that polynomial rings inherit this good behavior, laying the foundation for
modern algebraic geometry.

Today, commutative algebra---the study of commutative rings and their modules---forms
the technical core of algebraic geometry, algebraic number theory, and parts of
combinatorics. The language of Spec, localization, and sheaves connects pure algebra
to geometry in a way that has been decisive since Grothendieck's reformulation in the
1950s and 1960s.

This paper covers the following topics:
\begin{enumerate}[label=(\roman*)]
  \item Basic definitions and examples of rings and fields;
  \item Ideals, quotient rings, and the Chinese remainder theorem;
  \item Ring homomorphisms and the isomorphism theorems;
  \item Polynomial rings, Gauss's lemma, and Eisenstein's criterion;
  \item The hierarchy UFD $\supset$ PID $\supset$ Euclidean domain;
  \item Noetherian and Artinian rings, the Hilbert basis theorem, Nakayama's lemma;
  \item Localization and Dedekind rings;
  \item Gröbner bases and the Buchberger algorithm.
\end{enumerate}

\medskip

\noindent\textbf{Notation.} Throughout, all rings are assumed to have a multiplicative
identity $1$, and ring homomorphisms are required to map $1$ to $1$.
We write $\mathbb{Z}$, $\mathbb{Q}$, $\mathbb{R}$, $\mathbb{C}$ for the integers,
rationals, reals, and complex numbers, and $\mathbb{F}_p = \mathbb{Z}/p\mathbb{Z}$
for the field with $p$ elements.

% ============================================================
\section{Ring Basics}
% ============================================================

\begin{definition}[Ring]
A \emph{ring} is a set $R$ together with two binary operations $+$ (addition) and
$\cdot$ (multiplication) satisfying:
\begin{enumerate}[label=(\roman*)]
  \item $(R,+)$ is an abelian group with identity element $0$;
  \item $(R,\cdot)$ is a monoid with identity element $1$;
  \item Multiplication distributes over addition: $a(b+c) = ab + ac$ and
        $(a+b)c = ac + bc$ for all $a,b,c \in R$.
\end{enumerate}
The ring is \emph{commutative} if $ab = ba$ for all $a,b \in R$.
\end{definition}

\begin{example}\label{ex:rings}
The following are the most important examples:
\begin{enumerate}[label=(\alph*)]
  \item $\mathbb{Z}$, $\mathbb{Q}$, $\mathbb{R}$, $\mathbb{C}$ with ordinary
        addition and multiplication;
  \item $\mathbb{Z}/n\mathbb{Z}$ (integers modulo $n$) for any $n \geq 1$;
  \item The polynomial ring $R[x_1,\ldots,x_n]$ over any ring $R$;
  \item The matrix ring $M_n(R)$---noncommutative for $n \geq 2$;
  \item The ring of continuous functions $C([0,1])$ with pointwise operations;
  \item Quadratic integer rings $\mathbb{Z}[\sqrt{d}] = \{a + b\sqrt{d} \mid a,b \in \mathbb{Z}\}$
        for squarefree $d \in \mathbb{Z}$.
\end{enumerate}
\end{example}

\begin{definition}[Special ring classes]
Let $R$ be a commutative ring with $1 \neq 0$.
\begin{itemize}
  \item $R$ is a \emph{domain} (or \emph{integral domain}) if $ab = 0$ implies
        $a = 0$ or $b = 0$.
  \item $R$ is a \emph{field} if every nonzero element has a multiplicative inverse.
  \item The \emph{group of units} is $R^\times = \{u \in R \mid \exists v: uv = 1\}$.
  \item The \emph{characteristic} $\chr(R)$ is the smallest positive integer $n$ with
        $n \cdot 1 = 0$, or $0$ if no such $n$ exists.
\end{itemize}
\end{definition}

\begin{proposition}\label{prop:field-domain}
Every field is an integral domain. Conversely, every finite integral domain is a field.
\end{proposition}

\begin{proof}
If $F$ is a field and $ab = 0$ with $a \neq 0$, then $b = a^{-1}(ab) = 0$.
For the converse, let $R$ be a finite domain and $a \in R \setminus \{0\}$.
The map $x \mapsto ax$ is injective (as $R$ is a domain) and the domain $R$ is
finite, so this map is surjective. In particular, $ax = 1$ has a solution.
\end{proof}

\begin{example}[Units and zero divisors]
In $\mathbb{Z}/12\mathbb{Z}$: the units are $\{1,5,7,11\}$ (elements coprime to
$12$), and $2 \cdot 6 = 0$ so $2$ is a zero divisor.
In $\mathbb{Z}[i]$ (Gaussian integers): $(\mathbb{Z}[i])^\times = \{\pm 1, \pm i\}$.
\end{example}

\begin{theorem}[Wedderburn's little theorem]
Every finite division ring is commutative (hence a field).
\end{theorem}

\begin{proof}[Proof sketch]
One uses the class equation for the multiplicative group and properties of
cyclotomic polynomials to show that the center of the division ring must
equal the whole ring. See \cite[Thm.~13.1]{Lang2002}.
\end{proof}

% ============================================================
\section{Ideals and Quotient Rings}
% ============================================================

\begin{definition}[Ideal]
A subset $I \subseteq R$ is a \emph{(two-sided) ideal} if:
\begin{enumerate}[label=(\roman*)]
  \item $(I,+)$ is a subgroup of $(R,+)$;
  \item $rI \subseteq I$ and $Ir \subseteq I$ for all $r \in R$
        (in the commutative case these two conditions coincide).
\end{enumerate}
For $a_1,\ldots,a_n \in R$, the ideal \emph{generated} by these elements is
\[
  (a_1,\ldots,a_n) = \Bigl\{\sum_{i=1}^{n} r_i a_i \,\Big|\, r_i \in R\Bigr\}.
\]
An ideal of the form $(a) = aR$ is called a \emph{principal ideal}.
\end{definition}

\begin{definition}[Prime and maximal ideals]
Let $R$ be a commutative ring and $\mathfrak{p}, \mathfrak{m} \subsetneq R$ ideals.
\begin{itemize}
  \item $\mathfrak{p}$ is a \emph{prime ideal} if $ab \in \mathfrak{p}$ implies
        $a \in \mathfrak{p}$ or $b \in \mathfrak{p}$.
  \item $\mathfrak{m}$ is a \emph{maximal ideal} if there is no ideal $I$ with
        $\mathfrak{m} \subsetneq I \subsetneq R$.
\end{itemize}
\end{definition}

\begin{theorem}[Characterization via quotient rings]
Let $I$ be an ideal in a commutative ring $R$.
\begin{enumerate}[label=(\roman*)]
  \item $I$ is prime if and only if $R/I$ is an integral domain.
  \item $I$ is maximal if and only if $R/I$ is a field.
  \item Every maximal ideal is prime (since every field is an integral domain).
\end{enumerate}
\end{theorem}

\begin{proof}
(i) $ab + I = (a+I)(b+I) = 0$ in $R/I$ iff $ab \in I$. The domain condition in $R/I$
says exactly that this implies $a \in I$ or $b \in I$.
(ii) $R/I$ is a field iff every nonzero element $a + I$ is a unit, i.e., there exists
$b$ with $ab - 1 \in I$. This is equivalent to $I$ being maximal: for any $r \notin I$,
the ideal $(I,r) = R$ gives $ar + j = 1$ for some $a \in R, j \in I$.
\end{proof}

\begin{example}\label{ex:ideals-Z}
In $\mathbb{Z}$: the prime ideals are $(0)$ and $(p)$ for primes $p$.
The maximal ideals are exactly the $(p)$. The ideal $(6) = (2)(3)$ is neither prime
nor maximal: $\mathbb{Z}/(6) \cong \mathbb{Z}/2 \times \mathbb{Z}/3$ is not a domain.
\end{example}

\begin{theorem}[Chinese Remainder Theorem]\label{thm:crt}
Let $R$ be a commutative ring and $I_1,\ldots,I_k$ pairwise coprime ideals
(meaning $I_i + I_j = R$ for $i \neq j$). Then the natural map
\[
  R \,\Big/\, \bigcap_{i=1}^k I_i \;\xrightarrow{\;\sim\;}\; \prod_{i=1}^k R/I_i
\]
is a ring isomorphism.
\end{theorem}

\begin{proof}
The map $\varphi\colon R \to \prod R/I_i$ given by $r \mapsto (r + I_i)_i$ is a
ring homomorphism with kernel $\bigcap I_i$.
Surjectivity (and hence bijectivity after passing to the quotient) follows from the
coprimality assumption by an inductive argument: for each $j$, there exist $e_j \in R$
with $e_j \equiv 1 \pmod{I_j}$ and $e_j \equiv 0 \pmod{I_i}$ for $i \neq j$.
These elements $e_j$ are the ``idempotents'' that allow us to hit any target
$(a_1,\ldots,a_k)$ by taking $r = \sum_j a_j e_j$.
\end{proof}

\begin{corollary}
$\mathbb{Z}/n\mathbb{Z} \cong \prod_{p^k \| n} \mathbb{Z}/p^k\mathbb{Z}$
(product over prime power factors of $n$).
\end{corollary}

The \emph{prime spectrum} $\Spec(R)$ is the set of all prime ideals of $R$,
endowed with the \emph{Zariski topology}: closed sets are of the form
$V(I) = \{\mathfrak{p} \in \Spec(R) \mid I \subseteq \mathfrak{p}\}$.
For $R = \mathbb{Z}/n\mathbb{Z}$, the prime ideals correspond exactly to the
prime divisors of $n$ (plus the zero ideal when $n = p$ is prime).

% ============================================================
\section{Ring Homomorphisms and Isomorphism Theorems}
% ============================================================

\begin{definition}[Ring homomorphism]
A map $\varphi\colon R \to S$ between rings is a \emph{ring homomorphism} if
\[
  \varphi(a+b) = \varphi(a) + \varphi(b),\quad
  \varphi(ab) = \varphi(a)\varphi(b),\quad
  \varphi(1_R) = 1_S.
\]
The \emph{kernel} is $\ker\varphi = \{r \in R \mid \varphi(r) = 0\}$, an ideal of $R$.
The \emph{image} is $\operatorname{Im}(\varphi)$, a subring of $S$.
\end{definition}

\begin{theorem}[First Isomorphism Theorem for Rings]
Let $\varphi\colon R \to S$ be a ring homomorphism. Then
\[
  R/\ker\varphi \;\cong\; \operatorname{Im}(\varphi).
\]
\end{theorem}

\begin{proof}
The map $\bar\varphi\colon R/\ker\varphi \to \operatorname{Im}(\varphi)$ defined by
$\bar\varphi(r + \ker\varphi) = \varphi(r)$ is well defined (independent of
representative), a ring homomorphism, injective (since $\ker\bar\varphi = 0$),
and surjective onto $\operatorname{Im}(\varphi)$ by definition.
\end{proof}

\begin{theorem}[Second and Third Isomorphism Theorems]
Let $R$ be a ring, $I \subseteq J$ ideals of $R$, and $I, J$ ideals of $R$.
\begin{enumerate}[label=(\roman*)]
  \item \textbf{Second:} If $S$ is a subring and $I$ is an ideal of $R$, then
        $S + I$ is a subring, $S \cap I$ is an ideal of $S$, and
        $S/(S \cap I) \cong (S+I)/I$.
  \item \textbf{Third:} $J/I$ is an ideal of $R/I$ and
        $(R/I)/(J/I) \cong R/J$.
  \item \textbf{Correspondence theorem:} There is a bijection between
        ideals of $R/I$ and ideals of $R$ containing $I$.
\end{enumerate}
\end{theorem}

\begin{example}
The evaluation homomorphism $\varphi\colon \mathbb{R}[x] \to \mathbb{C}$ defined by
$\varphi(f) = f(i)$ has kernel $(x^2 + 1)$ and image $\mathbb{C}$, giving
$\mathbb{R}[x]/(x^2+1) \cong \mathbb{C}$.
\end{example}

\begin{definition}[Nilradical and Jacobson radical]
For a commutative ring $R$:
\begin{itemize}
  \item The \emph{nilradical} is $\Nil(R) = \{r \in R \mid r^n = 0 \text{ for some } n\} = \bigcap_{\mathfrak{p} \in \Spec(R)} \mathfrak{p}$.
  \item The \emph{Jacobson radical} is $\Jac(R) = \bigcap_{\mathfrak{m} \text{ maximal}} \mathfrak{m}$.
\end{itemize}
\end{definition}

% ============================================================
\section{Polynomial Rings}
% ============================================================

\begin{definition}[Polynomial ring]
For a commutative ring $R$, the \emph{polynomial ring} $R[x]$ consists of all formal
sums $f = \sum_{k=0}^n a_k x^k$ with $a_k \in R$, subject to the usual addition and
convolution multiplication. More generally, $R[x_1,\ldots,x_n]$ denotes polynomials
in $n$ variables.
\end{definition}

\begin{theorem}[Gauss's Lemma]
Let $R$ be a UFD with fraction field $K = \Frac(R)$.
A polynomial $f \in R[x]$ of degree $\geq 1$ is irreducible in $R[x]$ if and only if
$f$ is \emph{primitive} (the gcd of its coefficients is $1$) and irreducible in $K[x]$.
Consequently, $R[x]$ is a UFD whenever $R$ is.
\end{theorem}

\begin{proof}
The key step is the ``Gauss product formula'': if $f, g \in R[x]$ are primitive,
so is $fg$. This is proved by assuming otherwise and reducing modulo a prime divisor
of a coefficient of $fg$ to obtain a contradiction in $(R/(p))[x]$, which is a domain
when $p$ is prime. From this one deduces that content is multiplicative, and the
equivalence of irreducibility in $R[x]$ versus $K[x]$ follows.
See \cite[IV, \S2]{Lang2002}.
\end{proof}

\begin{theorem}[Eisenstein's Criterion]\label{thm:eisenstein}
Let $R$ be a UFD, $\mathfrak{p}$ a prime ideal, and
$f = a_n x^n + \cdots + a_1 x + a_0 \in R[x]$ with $n \geq 1$.
Suppose:
\begin{enumerate}[label=(\roman*)]
  \item $a_n \notin \mathfrak{p}$,
  \item $a_0, a_1, \ldots, a_{n-1} \in \mathfrak{p}$,
  \item $a_0 \notin \mathfrak{p}^2$.
\end{enumerate}
Then $f$ is irreducible in $\Frac(R)[x]$.
\end{theorem}

\begin{proof}
Suppose $f = gh$ with $g, h \in R[x]$ of positive degree. Reducing modulo $\mathfrak{p}$
gives $\bar{f} = a_n x^n = \bar{g}\bar{h}$ in $(R/\mathfrak{p})[x]$, a domain.
So $\bar{g} = cx^r$ and $\bar{h} = dx^s$ for some $r + s = n$, $r, s \geq 1$.
Then $\mathfrak{p} \mid g(0)$ and $\mathfrak{p} \mid h(0)$, hence
$\mathfrak{p}^2 \mid g(0)h(0) = a_0$, contradicting (iii).
\end{proof}

\begin{example}
\begin{enumerate}[label=(\alph*)]
  \item $x^p - p$ is irreducible over $\mathbb{Q}$ by Eisenstein with $\mathfrak{p}=(p)$.
  \item The $p$-th cyclotomic polynomial $\Phi_p(x) = x^{p-1} + \cdots + x + 1$ is
        irreducible: substitute $x \mapsto x+1$ and apply Eisenstein with the prime $p$.
\end{enumerate}
\end{example}

\begin{theorem}
For a commutative ring $R$, the polynomial ring $R[x]$ satisfies:
\begin{itemize}
  \item If $R$ is a domain, so is $R[x]$.
  \item $\chr(R[x]) = \chr(R)$.
  \item $(R[x])^\times = R^\times$ (if $R$ is a domain).
  \item $\mathbb{F}_q[x]$ (for $q = p^n$ a prime power) is a Euclidean domain,
        hence a PID and a UFD.
\end{itemize}
\end{theorem}

% ============================================================
\section{Unique Factorization: UFDs, PIDs, and Euclidean Domains}
% ============================================================

\begin{definition}[Divisibility and irreducibles]
In an integral domain $R$:
\begin{itemize}
  \item $a$ \emph{divides} $b$ (written $a \mid b$) if $b = ac$ for some $c \in R$.
  \item $p \in R$ is \emph{irreducible} if $p = ab$ implies $a$ or $b$ is a unit.
  \item $p \in R$ is \emph{prime} if $(p)$ is a prime ideal, i.e., $p \mid ab$ implies
        $p \mid a$ or $p \mid b$.
  \item Every prime is irreducible; in a UFD, the converse holds.
\end{itemize}
\end{definition}

\begin{definition}[UFD, PID, Euclidean domain]
\begin{itemize}
  \item A domain $R$ is a \emph{unique factorization domain} (UFD) if every nonzero
        nonunit has a unique factorization into irreducibles (up to order and units).
  \item $R$ is a \emph{principal ideal domain} (PID) if every ideal is principal.
  \item $R$ is a \emph{Euclidean domain} if there exists a function
        $\delta\colon R \setminus \{0\} \to \mathbb{N}_0$ such that for all
        $a \in R$, $b \in R \setminus \{0\}$ there exist $q, r \in R$ with
        $a = qb + r$ and either $r = 0$ or $\delta(r) < \delta(b)$.
\end{itemize}
\end{definition}

\begin{theorem}[Hierarchy of ring classes]\label{thm:hierarchy}
The following implications hold and none is reversible:
\[
  \text{Euclidean domain} \;\Rightarrow\; \text{PID} \;\Rightarrow\; \text{UFD}
  \;\Rightarrow\; \text{Integral domain}.
\]
\end{theorem}

\begin{proof}[Proof sketch]
\textit{Euclidean $\Rightarrow$ PID:} Let $I \neq 0$ be an ideal and $b \in I$
with $\delta(b)$ minimal. Then for any $a \in I$ write $a = qb + r$. Since
$r = a - qb \in I$ and $\delta(b)$ is minimal, $r = 0$, so $a \in (b)$.

\textit{PID $\Rightarrow$ UFD:} In a PID, every irreducible is prime (since
$(p)$ is maximal, hence prime). Existence of factorizations follows from
the ascending chain condition (all PIDs are Noetherian). Uniqueness follows from
the prime property.

\textit{UFD $\Rightarrow$ Domain:} By definition.

Counterexamples to converses: $\mathbb{Z}[\sqrt{-5}]$ is a UFD? No: $6 = 2 \cdot 3 =
(1+\sqrt{-5})(1-\sqrt{-5})$---two distinct factorizations, so it is not a UFD.
The ring $\mathbb{Z}[x]$ is a UFD but not a PID (the ideal $(2,x)$ is not principal).
The ring $k[x,y]$ is a UFD but not a PID for any field $k$.
There exist PIDs that are not Euclidean (e.g., $\mathbb{Z}[(1+\sqrt{-19})/2]$).
\end{proof}

\begin{example}[Euclidean domains]
\begin{enumerate}[label=(\alph*)]
  \item $\mathbb{Z}$ with $\delta(n) = |n|$;
  \item $k[x]$ for any field $k$ with $\delta(f) = \deg(f)$;
  \item The Gaussian integers $\mathbb{Z}[i]$ with $\delta(a+bi) = a^2 + b^2$
        (the norm). The Euclidean algorithm allows us to write $\gcd(m,n)$ in
        $\mathbb{Z}[i]$ explicitly. Two-square theorem: $p = a^2 + b^2$ for a prime
        $p$ iff $p = 2$ or $p \equiv 1 \pmod 4$.
\end{enumerate}
\end{example}

\begin{theorem}[Unique factorization in $\mathbb{Z}[i]$]
The Gaussian integers form a Euclidean domain (hence a PID and UFD). The units are
$\{\pm 1, \pm i\}$. The Gaussian primes are:
\begin{itemize}
  \item $(1+i)$ (lying over $2$);
  \item $p \in \mathbb{Z}$ prime with $p \equiv 3 \pmod 4$ (remain prime in $\mathbb{Z}[i]$);
  \item $a + bi$ with $a^2 + b^2 = p$ for primes $p \equiv 1 \pmod 4$.
\end{itemize}
\end{theorem}

% ============================================================
\section{Noetherian and Artinian Rings}
% ============================================================

\begin{definition}[Chain conditions]
A commutative ring $R$ is \emph{Noetherian} if it satisfies the \emph{ascending
chain condition} (ACC) on ideals: every chain $I_1 \subseteq I_2 \subseteq \cdots$
of ideals eventually stabilizes ($I_n = I_{n+1} = \cdots$ for some $n$).

Equivalently, $R$ is Noetherian iff every ideal is finitely generated.

$R$ is \emph{Artinian} if it satisfies the \emph{descending chain condition} (DCC):
every chain $I_1 \supseteq I_2 \supseteq \cdots$ stabilizes.
\end{definition}

\begin{theorem}[Hilbert Basis Theorem]\label{thm:hilbert-basis}
If $R$ is a Noetherian ring, then the polynomial ring $R[x]$ is also Noetherian.
\end{theorem}

\begin{proof}
Let $I \subseteq R[x]$ be an ideal. We show $I$ is finitely generated.
For each $n \geq 0$, let $J_n = \{a \in R \mid \text{there exists } f \in I
\text{ of degree } n \text{ with leading coeff.\ } a\} \cup \{0\}$.
Each $J_n$ is an ideal of $R$, and $J_0 \subseteq J_1 \subseteq J_2 \subseteq \cdots$
(multiplying by $x$ raises degree). Since $R$ is Noetherian, this chain stabilizes:
$J_N = J_{N+1} = \cdots$ for some $N$. Each $J_n$ is finitely generated: pick
generators $a_{n,1},\ldots,a_{n,k_n} \in J_n$ and corresponding polynomials
$f_{n,j} \in I$ of degree $n$ with leading coefficient $a_{n,j}$.
The finite set $\{f_{n,j}\}_{n \leq N, j \leq k_n}$ generates $I$ (by a
degree-induction argument). Hence $I$ is finitely generated.
\end{proof}

\begin{corollary}
For any field $k$, the polynomial ring $k[x_1,\ldots,x_n]$ is Noetherian.
Hence every ideal in $k[x_1,\ldots,x_n]$ has a finite generating set (a
\emph{Gröbner basis}, see Section~\ref{sec:groebner}).
\end{corollary}

\begin{theorem}[Nakayama's Lemma]\label{thm:nakayama}
Let $R$ be a local ring with maximal ideal $\mathfrak{m}$, and let $M$ be a
finitely generated $R$-module. If $M = \mathfrak{m}M$, then $M = 0$.

More generally: if $N \subseteq M$ is a submodule and $M = N + \mathfrak{m}M$,
then $M = N$.
\end{theorem}

\begin{proof}
Let $m_1,\ldots,m_k$ generate $M$. Since $M = \mathfrak{m}M$, we can write
$m_i = \sum_{j=1}^k a_{ij} m_j$ with $a_{ij} \in \mathfrak{m}$.
In matrix form, $(I - A)\mathbf{m} = 0$ where $A = (a_{ij})$.
The matrix $I - A$ has entries $\delta_{ij} - a_{ij}$; since $a_{ij} \in \mathfrak{m}$,
the diagonal entries $1 - a_{ii}$ are units (the only nonunits in a local ring lie
in $\mathfrak{m}$). Hence $\det(I-A)$ is a unit, so $I-A$ is invertible, giving $\mathbf{m} = 0$.
\end{proof}

\begin{corollary}
Let $R$ be local with maximal ideal $\mathfrak{m}$ and $M$ finitely generated.
Then $m_1,\ldots,m_k \in M$ generate $M$ iff their images in $M/\mathfrak{m}M$ generate
$M/\mathfrak{m}M$ as an $R/\mathfrak{m}$-vector space.
\end{corollary}

\begin{theorem}[Hopkins--Levitzki]
A Noetherian Artinian ring $R$ has a composition series as a module over itself.
Moreover: $R$ is Artinian if and only if $R$ is Noetherian of Krull dimension $0$.
\end{theorem}

The \emph{Krull dimension} $\dim(R)$ of a ring $R$ is the supremum of lengths of
chains of prime ideals $\mathfrak{p}_0 \subsetneq \mathfrak{p}_1 \subsetneq \cdots \subsetneq \mathfrak{p}_n$.
For fields, $\dim = 0$; for $\mathbb{Z}$ and $k[x]$, $\dim = 1$; for $k[x_1,\ldots,x_n]$, $\dim = n$.

% ============================================================
\section{Localization}
% ============================================================

\begin{definition}[Localization]
Let $R$ be a commutative ring and $S \subseteq R$ a \emph{multiplicative set}
(closed under multiplication, containing $1$). The \emph{localization}
$S^{-1}R$ is the set of equivalence classes of pairs $(a, s) \in R \times S$,
where $(a,s) \sim (a',s')$ iff $\exists t \in S: t(as' - a's) = 0$, with operations
\[
  \frac{a}{s} + \frac{a'}{s'} = \frac{as' + a's}{ss'}, \qquad
  \frac{a}{s} \cdot \frac{a'}{s'} = \frac{aa'}{ss'}.
\]
There is a natural ring homomorphism $\iota\colon R \to S^{-1}R$, $r \mapsto r/1$.
\end{definition}

\begin{example}[Key instances]
\begin{enumerate}[label=(\alph*)]
  \item $S = R \setminus \{0\}$: $S^{-1}R = \Frac(R)$ is the fraction field of a domain.
  \item $S = R \setminus \mathfrak{p}$ for a prime ideal $\mathfrak{p}$:
        $R_\mathfrak{p} = S^{-1}R$ is the \emph{localization at} $\mathfrak{p}$,
        a local ring with maximal ideal $\mathfrak{p}R_\mathfrak{p}$.
  \item $S = \{1, f, f^2, \ldots\}$ for $f \in R$: $R_f = S^{-1}R$ inverts $f$.
\end{enumerate}
\end{example}

\begin{theorem}[Universal property of localization]
The map $\iota\colon R \to S^{-1}R$ is initial among ring homomorphisms $\varphi\colon R \to T$
such that $\varphi(s)$ is a unit in $T$ for all $s \in S$. That is, any such $\varphi$
factors uniquely through $\iota$.
\end{theorem}

\begin{theorem}[Ideals in localizations]
There is a bijection between:
\begin{itemize}
  \item Prime ideals of $S^{-1}R$, and
  \item Prime ideals $\mathfrak{p}$ of $R$ with $\mathfrak{p} \cap S = \emptyset$.
\end{itemize}
Under this bijection, $\mathfrak{p} \leftrightarrow S^{-1}\mathfrak{p}$.
\end{theorem}

\begin{definition}[Dedekind domain]
An integral domain $R$ is a \emph{Dedekind domain} if:
\begin{enumerate}[label=(\roman*)]
  \item $R$ is Noetherian,
  \item $R$ is integrally closed in $\Frac(R)$,
  \item Every nonzero prime ideal of $R$ is maximal (i.e., $\dim R \leq 1$).
\end{enumerate}
\end{definition}

\begin{theorem}[Unique ideal factorization in Dedekind domains]
In a Dedekind domain $R$, every nonzero ideal factors uniquely (up to order) into
prime ideals:
\[
  I = \mathfrak{p}_1^{e_1} \cdots \mathfrak{p}_r^{e_r}.
\]
\end{theorem}

\begin{proof}[Proof sketch]
The key steps are: (1) every ideal contains a product of prime ideals (using
Noetherian induction); (2) for each prime $\mathfrak{p}$, the localization $R_\mathfrak{p}$
is a DVR (discrete valuation ring), i.e., a PID with a unique maximal ideal; and
(3) the factorization is assembled from the local data. See \cite{Neukirch1999}.
\end{proof}

\begin{example}[Rings of integers as Dedekind domains]
The ring of integers $\mathcal{O}_K$ of a number field $K = \mathbb{Q}(\sqrt{d})$ is
a Dedekind domain. Its integral closure in $K$ is:
\[
  \mathcal{O}_K = \begin{cases}
    \mathbb{Z}\!\left[\tfrac{1+\sqrt{d}}{2}\right] & \text{if } d \equiv 1 \pmod 4, \\
    \mathbb{Z}[\sqrt{d}] & \text{if } d \equiv 2 \text{ or } 3 \pmod 4.
  \end{cases}
\]
The \emph{class group} $\Cl(\mathcal{O}_K)$ measures the failure of $\mathcal{O}_K$
to be a PID. One has $h(d) = |\Cl(\mathcal{O}_K)| = 1$ iff $\mathcal{O}_K$ is a PID.
The Stark--Heegner theorem classifies all imaginary quadratic fields with $h = 1$:
$d \in \{-1, -2, -3, -7, -11, -19, -43, -67, -163\}$.
\end{example}

% ============================================================
\section{Noether Normalization and Integral Extensions}
% ============================================================

\begin{definition}[Integral element]
Let $R \subseteq S$ be a ring extension. An element $s \in S$ is \emph{integral over $R$}
if it satisfies a monic polynomial equation $s^n + a_{n-1}s^{n-1} + \cdots + a_0 = 0$
with $a_i \in R$. The \emph{integral closure} of $R$ in $S$ is the subring
$\bar{R} = \{s \in S \mid s \text{ integral over } R\}$.
$R$ is \emph{integrally closed} if $\bar{R} = R$ (within $\Frac(R)$).
\end{definition}

\begin{theorem}[Noether Normalization Lemma]
Let $k$ be a field and $A = k[x_1,\ldots,x_n]/I$ a finitely generated $k$-algebra
($I$ an ideal). Then there exist algebraically independent elements
$y_1,\ldots,y_d \in A$ such that $A$ is a finite (hence integral) extension of
$k[y_1,\ldots,y_d]$. Here $d = \dim A$ is the Krull dimension of $A$.
\end{theorem}

\begin{proof}[Proof sketch]
By induction on $n$. If $I = 0$ then $A = k[x_1,\ldots,x_n]$ and $y_i = x_i$ works.
Otherwise pick $f \in I$, $f \neq 0$. Make a change of variables $x_i' = x_i - x_1^{c^{i-1}}$
(for suitable $c$) to ensure $f$ becomes monic in $x_1$ after this substitution.
Then $x_1$ is integral over the smaller ring $k[x_2',\ldots,x_n']$, and one proceeds
by induction. See \cite[Thm.~5.10]{AtiyahMacdonald1969}.
\end{proof}

\begin{corollary}
For an affine variety $V \subseteq \mathbb{A}^n_k$ defined by an ideal $I \subseteq k[x_1,\ldots,x_n]$:
\[
  \dim V = \dim k[V] = \dim k[x_1,\ldots,x_n]/I.
\]
In particular, $\dim \mathbb{A}^n_k = n$ and $\dim$ of a hyperfläche $V(f) = n-1$.
\end{corollary}

% ============================================================
\section{Gröbner Bases}\label{sec:groebner}
% ============================================================

Gröbner bases, introduced by Bruno Buchberger in 1965~\cite{Buchberger1970},
provide a canonical set of generators for polynomial ideals and transform many
ideal-membership and algebraic-variety problems into algorithmic computations.

\begin{definition}[Monomial order]
A \emph{monomial order} on $k[x_1,\ldots,x_n]$ is a total ordering $\prec$ on
monomials $x^\alpha = x_1^{\alpha_1}\cdots x_n^{\alpha_n}$ that is compatible with
multiplication and has $1 \prec x^\alpha$ for all $\alpha \neq 0$.

Standard examples: \emph{lexicographic} (lex), \emph{graded reverse lexicographic} (grevlex),
and \emph{graded lexicographic} (grlex).
\end{definition}

\begin{definition}[Division algorithm and Gröbner basis]
Fix a monomial order $\prec$. For $f \in k[x_1,\ldots,x_n]$, write $\operatorname{LT}(f)$
for the \emph{leading term} (highest monomial times its coefficient).
A set $G = \{g_1,\ldots,g_s\} \subseteq I$ is a \emph{Gröbner basis} for the ideal
$I$ if
\[
  \langle \operatorname{LT}(g_1), \ldots, \operatorname{LT}(g_s) \rangle
  = \langle \operatorname{LT}(f) \mid f \in I \rangle.
\]
Equivalently, every element of $I$ reduces to $0$ modulo $G$.
\end{definition}

\begin{theorem}[Buchberger's Criterion]
$G = \{g_1,\ldots,g_s\}$ is a Gröbner basis for $I = (G)$ iff for all $i \neq j$,
the \emph{S-polynomial}
\[
  S(g_i, g_j) = \frac{\operatorname{lcm}(\operatorname{LM}(g_i), \operatorname{LM}(g_j))}{\operatorname{LT}(g_i)} g_i
  - \frac{\operatorname{lcm}(\operatorname{LM}(g_i), \operatorname{LM}(g_j))}{\operatorname{LT}(g_j)} g_j
\]
reduces to $0$ modulo $G$.
\end{theorem}

\begin{algorithm}[Buchberger's Algorithm]
\textbf{Input:} $F = \{f_1,\ldots,f_s\}$; \textbf{Output:} Gröbner basis $G$ for $(F)$.
\begin{enumerate}
  \item Set $G := F$.
  \item For all pairs $\{g_i, g_j\} \subseteq G$, compute $r = \overline{S(g_i,g_j)}^G$
        (remainder on division by $G$).
  \item If $r \neq 0$, add $r$ to $G$ and repeat.
  \item Return $G$.
\end{enumerate}
The algorithm terminates because the leading term ideal strictly increases at each step,
and $k[x_1,\ldots,x_n]$ is Noetherian.
\end{algorithm}

\begin{theorem}[Applications of Gröbner bases]
Given a Gröbner basis $G$ for $I$:
\begin{enumerate}[label=(\roman*)]
  \item \textbf{Ideal membership:} $f \in I$ iff $\overline{f}^G = 0$.
  \item \textbf{Solving systems:} Eliminate variables by computing Gröbner bases
        with lex order; the last variable satisfies a univariate polynomial.
  \item \textbf{Dimension:} $\dim k[x_1,\ldots,x_n]/I$ equals the Krull dimension
        computed from the leading monomial ideal.
  \item \textbf{Hilbert function:} The number of monomials \emph{not} in
        $\langle \operatorname{LT}(I) \rangle$ of degree $\leq d$ gives the
        Hilbert function of $k[x_1,\ldots,x_n]/I$.
\end{enumerate}
\end{theorem}

\begin{example}
Consider $I = (x^2 + y - 1,\; y^2 - x) \subseteq \mathbb{Q}[x,y]$ with lex order $x \succ y$.
Computing the Gröbner basis:
\[
  G = \{y^4 - y^2 + y - 1 + y^2,\; x - y^2\} \quad\text{(schematically)}.
\]
The lex basis exhibits $y$ satisfying a univariate polynomial, allowing solution
of the system $\{x^2 + y = 1, y^2 = x\}$ by back-substitution.
\end{example}

% ============================================================
\section{References}
% ============================================================

\begin{thebibliography}{10}

\bibitem{AtiyahMacdonald1969}
M.~F. Atiyah and I.~G. Macdonald.
\newblock \textit{Introduction to Commutative Algebra}.
\newblock Addison-Wesley, Reading, MA, 1969.
\newblock The classic concise introduction; covers localization, Noetherian rings, and
dimension theory.

\bibitem{Lang2002}
S.~Lang.
\newblock \textit{Algebra}, revised 3rd ed.
\newblock Springer, New York, 2002.
\newblock Comprehensive reference; Gauss's lemma, Eisenstein criterion, UFDs, and
field theory.

\bibitem{Matsumura1989}
H.~Matsumura.
\newblock \textit{Commutative Ring Theory}.
\newblock Cambridge Studies in Advanced Mathematics, vol.~8.
\newblock Cambridge University Press, Cambridge, 1989.
\newblock Advanced treatment of Noetherian/Artinian rings, dimension, flatness.

\bibitem{Neukirch1999}
J.~Neukirch.
\newblock \textit{Algebraic Number Theory}.
\newblock Grundlehren der mathematischen Wissenschaften, vol.~322.
\newblock Springer, Berlin, 1999.
\newblock Dedekind domains, ideal factorization, class groups from a number-theoretic
perspective.

\bibitem{Eisenbud1995}
D.~Eisenbud.
\newblock \textit{Commutative Algebra with a View toward Algebraic Geometry}.
\newblock Graduate Texts in Mathematics, vol.~150.
\newblock Springer, New York, 1995.
\newblock Thorough modern treatment connecting commutative algebra and geometry;
Gröbner bases in Chapter~15.

\bibitem{CoxLittleOShea2015}
D.~Cox, J.~Little, and D.~O'Shea.
\newblock \textit{Ideals, Varieties, and Algorithms}, 4th ed.
\newblock Undergraduate Texts in Mathematics.
\newblock Springer, New York, 2015.
\newblock Accessible introduction to Gröbner bases, Buchberger's algorithm, and
computational algebraic geometry.

\bibitem{Buchberger1970}
B.~Buchberger.
\newblock Ein algorithmisches Kriterium für die Lösbarkeit eines algebraischen
Gleichungssystems.
\newblock \textit{Aequationes Mathematicae}, 4(3):374--383, 1970.
\newblock Original paper introducing the Buchberger algorithm.

\bibitem{Noether1921}
E.~Noether.
\newblock Idealtheorie in Ringbereichen.
\newblock \textit{Mathematische Annalen}, 83(1--2):24--66, 1921.
\newblock Foundational paper establishing the axiomatic theory of ideals and the
ascending chain condition.

\bibitem{Hungerford1980}
T.~W. Hungerford.
\newblock \textit{Algebra}.
\newblock Graduate Texts in Mathematics, vol.~73.
\newblock Springer, New York, 1980.
\newblock Good reference for isomorphism theorems, polynomial rings, and modules.

\bibitem{Dummit2004}
D.~S. Dummit and R.~M. Foote.
\newblock \textit{Abstract Algebra}, 3rd ed.
\newblock Wiley, Hoboken, NJ, 2004.
\newblock Standard graduate textbook; comprehensive coverage of ring theory including
Noetherian rings and UFDs.

\end{thebibliography}

\end{document}
